Наверное, одна из самых популярных и полезных задач, которая связана с графами -- это нахождение пути из точки $A$ в точку $B$. В частности такую задачу решает, очевидно, навигатор ведь карта Москвы достаточно просто переносится в графовую интерпретацию и запрос <<построй маршрут от работы до дома>> сводится к задаче нахождения некоторой последовательности рёбер, по которым мы доезжаем до нашего уютного очага. Формально поставим нашу задачу.

\Def Пусть дан граф $G \langle V, E \rangle$. Путем из $u_a$ в $u_b$ будем называть конечную последовательность рёбер $P = \{(u_a, v_1), (v_1, v_2) \dots (v_n, u_b) \}$, что тождественно обозначению $P = \{ u_a \to v_1 \to \dots \to u_b \}$.

\Def Пусть $P = \{ u_a \to v_1 \to \dots u_b \}$ - путь. Тогда $|P| \overset{def}{=} n - 1$.

\Task Пусть дан граф $G \langle V, E \rangle$, а также $u_a, u_b \in V$. Необходимо найти путь из $u_a$ в $u_b$.

Но такую задачу мы в принципе уже научились решать, просто запуская DFS или BFS. Скорее, как пользователю навигатора, нам очень хотелось бы побыстрее добраться к семье, поэтому мы хотим получить \textit{наикратчайший путь}. Для того чтобы говорить о <<наикратчайшести>> нужно придать ребрам длину или же задать \textit{весовую функцию}.

\Def Пусть дан граф $G \langle V, E \rangle$. Весовой функцией будем называть отображение $w: E \to \R$.

Например, в навигаторе в качестве весовой функции можно считать некоторую комбинацию расстояния и, например, пробок.

\Def Пусть дан граф $G \langle V, E \rangle$, весовая функция $w: E \to \mathbb{R}$, а также $u, v \in V$ и некоторый путь $P = \{e_k\}_{k=1}^n$ от $u$ до $v$. Тогда весом пути примем: 
\begin{equation*}
    w(P) \overset{def}{=} \sum\limits_{k = 1}^n w(e_k)
\end{equation*}

\Def Пусть дан граф $G \langle V, E \rangle$, весовая функция $w: E \to \mathbb{R}$, а также $u, v \in V$. Кратчайшим путём из $u$ в $v$ будем называть путь минимальный по весу и его вес будем обозначать $\delta(u, v)$.

Собственно, вот задача, которую мы будем решать в данном разделе.

\Task Пусть дан граф $G \langle V, E \rangle$, весовая функция $w: E \to \R$, а также $u_a, u_b \in V$. Необходимо найти кратчайший путь из $u_a$ в $u_b$.

Всюду далее считаем, что задан некоторый граф $G\langle V, E \rangle$ и весовая функция $w: E \to \R$.

\Def Множество соседей вершины $u$ в графе будем обозначать $\alpha(u)$.
